\chapter{Theory}
\section{Mô hình Đối tượng Tính toán COKB}

Mô hình \textbf{COKB (Computational Objects Knowledge Base)} [1] là sự mở rộng của các hệ thống logic truyền thống, tích hợp khả năng tính toán mạnh mẽ vào các cấu trúc đối tượng, phục vụ như là định dạng lưu trữ và biểu diễn bộ não cốt lõi của hệ thống KBMS.

\subsection{Thành phần Hệ thống tri thức COKB}

Một cơ sở tri thức COKB được xác định bởi bộ 6 thành phần [1, 3]:
\begin{equation}
COKB = (C, H, R, Ops, Funcs, Rules)
\end{equation}

Trong đó:
*   \textbf{C (Concepts)}: Tập hợp các khái niệm hoặc lớp đối tượng tính toán.
*   \textbf{H (Hierarchy)}: Các quan hệ phân cấp đặc biệt hóa giữa các khái niệm (quan hệ IS-A).
*   \textbf{R (Relations)}: Tập các quan hệ ngữ nghĩa giữa các lớp đối tượng (ví dụ: song song, vuông góc).
*   \textbf{Ops (Operators)}: Các toán tử tính toán trên các miền giá trị (Số thực, Vector, Ma trận).
*   \textbf{Funcs (Functions)}: Các hàm xác định ánh xạ giữa các thuộc tính.
*   \textbf{Rules (Rules)}: Tập hợp các luật dẫn để suy diễn ra tri thức mới.

\subsection{Mô hình Đối tượng Tính toán (Computational Object)}

Mỗi thực thể (Object) trong hệ thống được biểu diễn bởi bộ ba thành phần [2]:
\begin{equation}
O = (Attrs, Facts, Rules)
\end{equation}

1.  \textbf{Attrs (Attributes)}: Tập các thuộc tính của đối tượng. Mỗi thuộc tính bản thân nó cũng có thể là một đối tượng tính toán khác (cấu trúc đệ quy) [4].
2.  \textbf{Facts (Facts)}: Các giá trị hoặc tính chất vốn có của đối tượng đã được xác định.
3.  \textbf{Rules (Internal Rules)}: Các quy tắc, phương trình nội tại ràng buộc mối quan hệ giữa các Attrs bên trong đối tượng đó.

\subsection{Phân cấp Khái niệm (Concept Levels)}

Trong mô hình COKB, các khái niệm được phân tầng dựa theo độ phức tạp của cấu trúc attributes:
*   \textbf{Cấp 0 (Base Concepts)}: Các kiểu dữ liệu cơ sở (Số thực - ℝ, Điểm - Point).
*   \textbf{Cấp 1}: Các khái niệm xây dựng trực tiếp từ cấp 0 (Đoạn thẳng, Góc).
*   \textbf{Cấp n}: Các khái niệm phức tạp hình thành từ các lớp thấp hơn (Tam giác, Tứ giác, Động cơ).

Việc phân cấp này giúp hệ thống quản lý tri thức theo hướng mô-đun hóa và hỗ trợ lan truyền kế thừa tri thức một cách tự động.

\section{Cơ chế Suy luận và Giải quyết vấn đề}

Dựa trên cấu trúc mô hình COKB, quá trình giải quyết bài toán thực chất là quá trình mở rộng tập sự thật dựa trên các quy tắc suy luận cưỡng bức, cho phép hệ thống "hiểu" và tự động phát sinh nội dung mà không cần người dùng nhập thủ công.

\subsection{Các Quy tắc Suy luận (Inference Rules)}

Hệ thống KBMS vận hành dựa trên 6 loại quy tắc suy luận chính (RC1 - RC6) [2]:

*   \textbf{RC1 (Vốn có)}: Sinh ra sự kiện từ thuộc tính mặc định của đối tượng (ví dụ: Tam giác cân có 2 góc đáy bằng nhau).
*   \textbf{RC2 (Mặc nhiên)}: Các phép toán gán, thay thế và bắc cầu cơ bản giữa các biến và hằng số.
*   \textbf{RC3 (Thay thế quan hệ)}: Sử dụng các quan hệ tính toán đã biết để tìm giá trị biến còn thiếu (Ohm's Law: $U, I \implies R$).
*   \textbf{RC4 (Luật dẫn)}: Kích hoạt các luật dạng `IF <Hypothesis> THEN <Goal>`.
*   \textbf{RC5 (Giải hệ phương trình)}: Kết hợp nhiều biến và phương trình để giải đồng thời bằng phương pháp Newton-Raphson hoặc Brent [2].
*   \textbf{RC6 (Hành vi nội bộ)}: Suy diễn dựa trên các đặc tính cấu thành (PART-OF) [2].

\subsection{Thuật toán Tìm bao đóng (F-Closure Algorithm)}

Đây là thuật toán trọng tâm của KBMS dùng để tìm tập đóng tri thức lớn nhất có thể suy ra được từ tập giả thiết \textbf{GT}.

\textbf{Mô tả thuật thuật giải}:
1.  Khởi tạo tập `KnownFacts = GT`.
2.  Lặp lại cho đến khi không còn sự thật mới (Fixed-point):
    *   Quét qua toàn bộ tập luật (Rules) và phương trình (Equations).
    *   Nếu điều kiện của luật thỏa mãn trong `KnownFacts`, đẩy kết luận vào tập đóng.
    *   Cập nhật `KnownFacts`.
3.  Kết quả: `FClosure(GT)` là tập hợp chứa toàn bộ tri thức tường minh và dẫn xuất.

\subsection{Chiến lược Giải bài toán: Heuristics}

Để tăng tốc độ suy diễn trong các không gian tri thức khổng lồ, KBMS sử dụng các quy tắc heuristics dựa trên lớp bài toán [2]:
*   \textbf{Heuristic 1}: Ưu tiên các luật có phần kết luận chứa biến mục tiêu (\textbf{Goal-directed}).
*   \textbf{Heuristic 2}: Ưu tiên giải quyết các đối tượng có mức độ xác định (Confidence Level) cao nhất trước.
*   \textbf{Heuristic 3}: Sử dụng đồ thị phụ thuộc biến để giới hạn phạm vi quét luật.

---

> [!NOTE]
> Việc tích hợp các luật heuristics giúp hệ thống giảm thiểu các bước suy diễn thừa, từ đó đạt được thời gian đáp ứng thời gian thực ngay cả trên các bộ dữ liệu lớn.

